% Generated draft: requires mathematical and source review.
\documentclass[11pt]{article}
\usepackage[margin=25mm]{geometry}
\usepackage{fontspec}
\setmainfont{Latin Modern Roman}
\usepackage{amsmath,amssymb,mathtools,unicode-math}
\setmathfont{Latin Modern Math}
\usepackage{enumitem,needspace,xurl,hyperref}
\setlength{\parindent}{0pt}
\setlength{\parskip}{5pt}
\setlength{\emergencystretch}{3em}
\newcommand{\N}{\mathbb{N}}
\newcommand{\Z}{\mathbb{Z}}
\newcommand{\Q}{\mathbb{Q}}
\newcommand{\R}{\mathbb{R}}
\newcommand{\C}{\mathbb{C}}
\newcommand{\F}{\mathbb{F}}
\newcommand{\A}{\mathbb{A}}
\newcommand{\PP}{\mathbb{P}}
\newcommand{\E}{\mathbb{E}}
\newcommand{\eps}{\varepsilon}
\newcommand{\dd}{\,\mathrm{d}}
\begin{document}
\section*{UnsolvedMath ID: 30003122}
\textbf{Source problem code:} catalogue entry
\subsection{Definitions and mathematical statement}
The supplied problem concerns a curve \(X/\Q\) of genus \(g\) and a nonconstant rational function \(t\in\Q(X)\) of degree \(\nu\geq2\). For an integral curve, the function-field interpretation of this degree is \(\nu=[\Q(X):\Q(t)]\). The source does not explicitly specify smoothness, projectivity, or geometric integrality, or its convention for genus when these hypotheses are absent.

Fix an algebraic closure \(\overline{\Q}\). For each positive integer \(n\), consider an algebraic point \(P_n\in X(\overline{\Q})\) at which \(t\) is defined and such that
\[
t(P_n)=n.
\]
Write \(\Q(P_n)\) for its field of definition over \(\Q\), viewed as a finite extension inside \(\overline{\Q}\). For a chosen sequence \((P_n)_{n\geq1}\), set
\[
D_N=\#\{\Q(P_n):1\leq n\leq N\},\qquad N\in\Z_{\geq1},
\]
where this displayed definition counts distinct subfields of the fixed algebraic closure.

The requested lower bound is
\[
D_N\geq c(g,\nu)N
\qquad\text{for every integer }N\geq N_0,
\]
with \(c(g,\nu)>0\), depending only on \(g\) and \(\nu\), and some integer \(N_0\geq1\).

The formulation is unresolved in several respects: the source does not say whether the bound must hold for every sequence of choices \(P_n\), or whether suitable choices may be made; it does not specify the permitted dependence of \(N_0\); and it does not explicitly distinguish equality of embedded fields from isomorphism over \(\Q\). If fields are to be counted up to \(\Q\)-isomorphism, \(D_N\) must instead denote the number of their isomorphism classes. The source also does not state hypotheses ensuring that every indicated fiber contains an eligible point on the chosen model of \(X\). These qualifications require clarification before the displayed target becomes a fully quantified conjecture.
\subsection{Short English statement}
The problem asks whether points above the integers under a rational function of degree at least two yield linearly many distinct number fields, with a positive constant depending only on the genus and degree. The curve hypotheses, choice of points, meaning of distinct fields, and dependence of the threshold require clarification.
\subsection{Sources}
\begin{itemize}
\item Supplied local UnsolvedMath registry record, permanent ID \(30003122\), title: \emph{Distinct Number Fields from Fibers of Rational Functions}; exact source problem code: \texttt{catalogue entry}. The record belongs to the Oberwolfach Reports Open Problems collection, gives proposed year 2016, and supplies no proposer attribution.
\item Supplied source citation: \emph{Diophantische Approximationen (2017)}, Oberwolfach Reports, DOI 10.4171/owr/2016/21, pp. 1099--1187. \url{https://doi.org/10.4171/owr/2016/21}. The supplied background describes Oberwolfach Reports as documenting mathematical workshops and their open-problem discussions.
\item Supplied literature source, labeled OWR report: \url{https://ems.press/content/serial-article-files/46625?nt=1}.
\item Supplied literature source, labeled Bilu--Luca: \url{https://arxiv.org/abs/1607.00904}.
\item Supplied literature source, labeled Hasson--Obus: \url{https://arxiv.org/abs/1706.05782}.
\item Supplied dated literature curation, checked \texttt{2026-08-21}: the record labels the problem partially solved and states that the weak diversity bound is proved up to a logarithmic loss and in several special cases, follows conditionally from the \(abc\) conjecture, and that the stated linear lower bound remains open unconditionally. This is the supplied dated assessment, not an independently verified literature conclusion. The background explicitly calls for independent verification of time-sensitive status and cited claims.
\item The supplied statement-status label is \texttt{corrected\_verified}; its accompanying explanation says that mathematical content was preserved while notation, typography, or omitted definitions were repaired for self-contained publication. This supplied editorial label does not establish mathematical correctness or resolve the formulation issues above.
\end{itemize}
\end{document}
